%======================================================================
% FILOSOFIA Y VISION - SANDRA DESKTOP CONTAINER
%======================================================================

\chapter*{Filosofia y Principios}
\addcontentsline{toc}{chapter}{Filosofia y Principios}
\markboth{Filosofia y Principios}{Filosofia y Principios}

\begin{quote}
\textit{``En un mundo de software efimero, Sandra Desktop Container establece un estandar de permanencia, seguridad y control.''}
\end{quote}

\section*{Vision}

\sdc{} no es simplemente un lanzador de aplicaciones; es un \textbf{Orquestador de Entornos Seguros}. Su proposito es abstraer la complejidad del sistema operativo subyacente (macOS, Linux, Windows) para ofrecer una interfaz unificada, segura y controlada donde las aplicaciones empresariales criticas pueden ejecutarse sin interferencias externas.

El futuro de SDC apunta hacia la \keyconcept{Computacion Descentralizada y Privada}, donde el contenedor gestiona no solo la ejecucion de la interfaz de usuario, sino tambien la identidad soberana del usuario, las llaves criptograficas y la persistencia de datos local-first, eliminando la dependencia absoluta de la nube para operaciones sensibles.

\section*{Principios Fundamentales}

\begin{quote}
\textbf{Seguridad por Diseno (Security by Design)}\\
\small
Cada componente de SDC ha sido concebido desde su origen con la seguridad como requisito fundamental, no como caracteristica anadida. Implementamos defensa en profundidad con multiples capas de validacion desde el kernel de Rust hasta la sanitizacion de la interfaz en Angular.
\end{quote}

\begin{enumerate}[leftmargin=*]
    \item \textbf{Seguridad de Memoria}: El backend escrito en Rust garantiza la ausencia de errores de segmentacion y condiciones de carrera, cumpliendo con los estandares mas altos de robustez.
    
    \item \textbf{Aislamiento (Sandboxing)}: Cada micro-aplicacion se ejecuta en un contexto controlado con politicas de seguridad de contenido estrictas, evitando ataques de Cross-Site Scripting (XSS) entre modulos.
    
    \item \textbf{Privacidad por Diseno}: Implementacion de cifrado estructural que garantiza que los datos sensibles solo sean legibles por el receptor autorizado, cumpliendo con GDPR y LOPD.
    
    \item \textbf{Principio de Minimo Privilegio}: Las micro-aplicaciones carecen de acceso al sistema de archivos o redes externas a menos que sea explicitamente autorizado mediante el puente IPC seguro.
\end{enumerate}

\section*{Arquitectura de Confianza}

La arquitectura de seguridad de SDC se organiza en capas:

\begin{table}[H]
\centering
\begin{tabular}{p{4cm}p{5cm}p{5cm}}
\toprule
\textbf{Capa} & \textbf{Componente} & \textbf{Funcion} \\
\midrule
Interfaz & Angular UI & Reactividad y sanitizacion \\
\midrule
Puente & Tauri IPC & Comunicacion segura \\
\midrule
Core & Rust & Seguridad de memoria \\
\midrule
Criptografia & AES-256-GCM & Cifrado autenticado \\
\midrule
Hashing & Argon2id & Proteccion de credenciales \\
\bottomrule
\end{tabular}
\caption{Componentes de la arquitectura de confianza}
\label{tab:arquitectura-confianza}
\end{table}

\section*{Compromiso con Estandares}

SDC ha sido disenado siguiendo rigurosamente principios que se alinean con normativas internacionales:

\begin{itemize}[leftmargin=*]
    \item \textbf{ISO/IEC 27001}: Controles de criptografia para asegurar confidencialidad e integridad
    \item \textbf{ISO/IEC 27032}: Directrices para proteccion de activos en el ciberespacio
    \item \textbf{COBIT 2019}: Marcos de gobernanza para objetivos de seguridad institucional
    \item \textbf{NIST Cybersecurity Framework}: Identificacion, proteccion, deteccion, respuesta y recuperacion
\end{itemize}

\section*{Bunker Digital}

El contenedor SDC actua como un \textbf{perimetro de seguridad} que rodea a la aplicacion. Cada accion esta vinculada a un SID (Secure ID) trazable, permitiendo auditorias forenses precisas sin comprometer la privacidad operativa.

\begin{quote}
\textbf{Responsabilidad Legal}\\
\small
SDC opera bajo un marco estricto de Hacking Etico y Legalidad Informatica. El Inspector de Red actua como herramienta de transparencia, permitiendo a los oficiales de seguridad auditar el trafico saliente sin recurrir a tecnicas de ``Man-in-the-Middle'' externas.
\end{quote}

\section*{Hacia el Futuro}

La hoja de ruta de SDC incluye:

\begin{enumerate}[leftmargin=*]
    \item \textbf{Identidad Soberana}: Gestion descentralizada de credenciales
    \item \textbf{Persistencia Local-First}: Datos criticos sin dependencia de nube
    \item \textbf{Orquestacion Distribuida}: Multiples nodos SDC sincronizados
    \item \textbf{Inteligencia Artificial Local}: Procesamiento de Sandra AI en el dispositivo
\end{enumerate}

\vspace{1cm}
\begin{center}
\fbox{\parbox{0.8\textwidth}{\centering
    \textcolor{sandraDark}{\large\bfseries ``Neutralidad del Hardware, Soberania del Usuario''}
}}
\end{center}
